\documentclass[11pt]{article}
\usepackage{preamble}
\setlength{\parskip}{4pt}

\begin{document}

\daybanner{AP Statistics \textperiodcentered\ Unit 2 \textperiodcentered\ Topic 2.10 \textperiodcentered\ B: 2026-11-05 \textperiodcentered\ E: 2026-11-25 \textperiodcentered\ TEACHER KEY}{Rare Is Not Never}

\sectionbanner{CED TETHER}

{\small
\begin{itemize}[leftmargin=1.5em,itemsep=2pt,topsep=2pt]
  \item \textbf{Skill 3.A} --- Determine relative frequencies, proportions, or probabilities using simulation or calculations.
  \item \textbf{EU UNC-3} --- Probabilistic reasoning allows us to anticipate patterns in data.
  \item \textbf{LO UNC-3.A} --- Estimate probabilities of binomial random variables using data from a simulation. [Skill 3.A]
  \item \textbf{EK UNC-3.A.1} --- A probability distribution can be constructed using the rules of probability or estimated with a simulation using random number generators.
  \item \textbf{EK UNC-3.A.2} --- A binomial random variable, X, counts the number of successes in n repeated independent trials, each trial having two possible outcomes (success or failure), with the probability of success p and the probability of failure 1 - p.
  \item \textbf{LO UNC-3.B} --- Calculate probabilities for a binomial distribution. [Skill 3.A]
  \item \textbf{EK UNC-3.B.1} --- The probability that a binomial random variable, X, has exactly x successes for n independent trials, when the probability of success is p, is calculated as P(X = x) = C(n, x) $\cdot$ p\textasciicircum{}x $\cdot$ (1 - p)\textasciicircum{}(n - x), for x = 0, 1, 2, \ldots{}, n. This is the binomial probability function.
\end{itemize}
}
\textbf{Essential question:} When does a low binomial probability count against a familiar chance model without proving it false?\par
\textbf{DOK-3 skill:} 4.B \textperiodcentered\ \textbf{FRQ pattern:} {\small\texttt{binomial-\allowbreak{}conditions-\allowbreak{}then-\allowbreak{}is-\allowbreak{}this-\allowbreak{}outcome-\allowbreak{}unusual}}\par
\textbf{DOK rationale:} Part (c) weighs competing claims using the day's numerical or graphical evidence and explains a limitation or consequence; the judgment requires linked reasoning beyond a calculation.\par

\frameworkphaseheader{Do Now (first take) $\rightarrow$ Explore (video + follow-along) $\rightarrow$ Exit (finish + turn in)}{1 $\rightarrow$ 3}{45}{%
\begin{itemize}[leftmargin=1.5em,itemsep=2pt,topsep=2pt]
  \item Show the board at the bell and collect a one-sentence first take without confirming a claim.
  \item Run the named video follow-along, then allow about 10 minutes for parts (a)--(c).
  \item Ask for evidence linked to a claim. Collect the sheet and hand-score part (c) with E/P/I.
\end{itemize}
}{%
\begin{itemize}[leftmargin=1.5em,itemsep=2pt,topsep=2pt]
  \item Commit to one claim or neither and cite the result or graph before the lesson.
  \item Complete the follow-along about BINS and exact versus cumulative binomial probability.
  \item Finish (a)--(c), revise the first take in the exit line, and turn in the sheet.
\end{itemize}
}{%
\begin{itemize}[leftmargin=1.5em,itemsep=2pt,topsep=2pt]
  \item Which number or visible feature supports your claim?
  \item Why does the other claim go beyond that evidence?
  \item What would you tell someone who chose the other claim?
\end{itemize}
}{Circulate and ask students to connect their choice to evidence. Score part (c) using the three required reasoning elements; do not require dropped extension work.}

\begin{callout}{\IconWarn\ WATCH FOR}{calloutred}
\begin{itemize}[leftmargin=1.5em,itemsep=2pt,topsep=2pt]
  \item A choice with copied numbers but no explanation of how they support it.
  \item Treating a model or average as a guaranteed individual outcome.
\end{itemize}
\end{callout}

\sectionbanner{SCORING PART (c) --- E / P / I}

\textbf{Expected elements}
\begin{itemize}[leftmargin=1.5em,itemsep=2pt,topsep=2pt]
  \item \textbf{reasoning-1} (required) --- Uses 0.0106 to distinguish unusual from impossible and assess both claims
  \item \textbf{reasoning-2} (required) --- Gives a plausible challenge to independence or constant p
  \item \textbf{reasoning-3} (required) --- Names comparable new shooting evidence and explains what pattern would matter
\end{itemize}

{\small\renewcommand{\arraystretch}{1.25}
\noindent\begin{tabular}{|p{0.5in}|p{5.9in}|}
\hline
\textbf{E} & Includes all three required reasoning elements above, with correct contextual evidence. \\ \hline
\textbf{P} & Includes at least one substantive correct reasoning link above, but omits or misstates another required element. \\ \hline
\textbf{I} & Includes no substantive correct reasoning link above; an unsupported choice or copied number alone is insufficient. \\ \hline
\end{tabular}}

\textbf{Common mistakes}
\begin{itemize}[leftmargin=1.5em,itemsep=2pt,topsep=2pt]
  \item Computing only $P(X=3)$ when the evidence is 3 or fewer
  \item Calling shots independent merely because there are 10 of them
  \item Treating a small probability as proof that the career probability changed
  \item Saying any result from 10 attempts must be ordinary because the sample is small
\end{itemize}

\clearpage
\focusbanner{TODAY'S PROBLEM --- annotated}

Suppose a basketball player has made 70 percent of her free throws over a long career. Tonight she makes 3 of 10 attempts. Let $X$ be the number she would make in 10 attempts if her chance on each shot were still 0.70. The graph shows the binomial model as expected frequencies per 1,000 sets of 10 shots.\par\smallskip \textbf{Claim A:} ``She is in a slump---making only 3 of 10 never happens at her career rate.''\par \textbf{Claim B:} ``With only 10 shots, 3 makes is ordinary chance variation.''\par Under this model, technology gives $P(X\leq3)=0.0106$ for 3 or fewer makes.\par
\begin{center}
\begin{tikzpicture}\begin{axis}[width=0.95\linewidth,height=1.8in,ybar,bar width=7pt,xmin=-0.6,xmax=10.6,ymin=0,ymax=290,xtick={0,1,2,3,4,5,6,7,8,9,10},ytick={0,100,200},xlabel={Made shots},ylabel={Count per 1,000},tick label style={font=\small},label style={font=\small},ymajorgrids]\addplot[fill=calloutblue,draw=dayblue] coordinates {(0,0.006) (1,0.138) (2,1.447) (3,9.002) (4,36.757) (5,102.919) (6,200.121) (7,266.828) (8,233.474) (9,121.061) (10,28.248)};\end{axis}\end{tikzpicture}
\end{center}
\begin{firsttakebox}{FIRST TAKE (5 min)}
Which claim, if either, do you trust more? Cite one detail from the result or graph.\par
\answer{Not graded. Preserve the possibility of neither claim; both statements intentionally overreach in different directions.}
\end{firsttakebox}

\textit{Rules callout ``BINS BEFORE BINOMIAL (keep on the page)'' is printed on the student sheet.}\par\medskip

\dokbadge{1}~\textbf{(a)} State what $X$ counts, and give $n$ and $p$.\par
\answer{$X$ counts makes in 10 attempts; $n=10$ and $p=0.70$.}

\dokbadge{2}~\textbf{(b)} Calculate $P(X=3)$ using the binomial formula. Explain how this event differs from $X\leq3$.\par
\answer{$P(X=3)=\binom{10}{3}(0.70)^3(0.30)^7=0.009001692\approx0.0090$. Exactly 3 includes one count; $X\leq3$ includes 0, 1, 2, and 3 makes.}

\dokbadge{3}~\textbf{(c)} Is either claim fully supported? Use the given chance of 3 or fewer makes to explain what each claim gets wrong. Name one reason independence or the same success chance might fail, and one kind of new evidence that would help. Write 3--4 sentences.\par
\answer{Neither claim is fully supported: $P(X\leq3)=0.0106$ makes this unusual under the model, contrary to B, but possible, contrary to A. One unusual set does not prove that the player's rate changed. Fatigue or injury could make a constant $0.70$ chance implausible. More attempts under comparable conditions would help: a sustained rate below $70\%$ would strengthen the slump explanation.}

\begin{summaryexitbox}{EXIT}
One thing the video changed about my first take: \hrulefill
\end{summaryexitbox}

\end{document}
